\documentclass[conference]{IEEEtran}
\IEEEoverridecommandlockouts
\usepackage{cite}
\usepackage{amsmath,amssymb,amsfonts}
\usepackage{graphicx}
\usepackage{booktabs}
\usepackage{url}
\usepackage{hyperref}
\usepackage{xcolor}
\def\BibTeX{{\rm B\kern-.05em{\sc i\kern-.025em b}\kern-.08em
    T\kern-.1667em\lower.7ex\hbox{E}\kern-.125emX}}

\begin{document}

\title{AgentShield: Provenance-Aware Defense Against\\Multilingual Prompt Injection in LLM Agents}

\author{\IEEEauthorblockN{[NAME]}
\IEEEauthorblockA{\textit{[DEPARTMENT]} \\
\textit{[INSTITUTION]}\\
[CITY, COUNTRY] \\
[EMAIL]}}

\maketitle

\begin{abstract}
LLM agents that call tools are vulnerable to \emph{indirect prompt injection}: instructions planted
in the content they read are executed as if issued by the user. Existing defenses are evaluated
almost entirely in English and are rarely tested against an adaptive attacker. We present
AgentShield, a layered, provenance-aware defense combining a de-obfuscating preprocessor, a
three-layer detector (multilingual rules, a fine-tuned multilingual classifier, and an LLM intent
check), canary tokens, and origin-based taint tracking, fused by logistic regression and enforced
by a policy firewall. We also build a multilingual indirect-injection dataset (English, Hindi,
Tamil, Hinglish, Tanglish) with benign hard negatives. On a public English injection test set, the
rule layer with learned fusion attains 0.78 precision at 0.23 recall and 3.6\% false positives at
$\sim$0.36\,ms per check, exceeding a keyword baseline, with the classifier layer carrying recall
once trained. Crucially, the provenance layers stop data exfiltration that content detection misses,
and the preprocessor neutralizes obfuscations that evade a keyword filter. We release the code,
dataset tooling, an SDK, a live dashboard, and a Docker deployment.
\end{abstract}

\begin{IEEEkeywords}
prompt injection, LLM agents, AI security, information-flow control, multilingual NLP
\end{IEEEkeywords}

\section{Introduction}
LLMs are increasingly deployed as \emph{agents} that read external content and invoke tools. The
same capability is the attack surface: an adversary who controls retrieved content can inject
instructions that hijack the agent, without ever contacting it directly~\cite{greshake}. As agents
gain the ability to send mail, move money, and call APIs, a hijack escalates from a wrong answer to
an irreversible action such as data exfiltration.

Two gaps motivate this work. \emph{Language}: agent-injection benchmarks~\cite{agentdojo,injecagent,bipia}
are English-only, yet low-resource and code-mixed inputs bypass safety training~\cite{multijail,lowres}.
\emph{Adaptivity}: adaptive attacks have broken eight published defenses, exceeding 50\% attack
success~\cite{adaptive}. AgentShield targets both and, critically, does not rely on detection alone:
provenance checks stop harmful actions even when detection is fooled.

\textbf{Contributions.} (i) The first multilingual indirect-injection dataset for agents with benign
hard negatives; (ii) a layered defense fusing probabilistic detection with deterministic provenance
(canary + taint); (iii) an adaptive-attacker evaluation and ablation; (iv) an open, reproducible
implementation.

\section{Related Work}
\emph{Prompt-level} defenses such as spotlighting~\cite{spotlight} mark untrusted text but rely on
model compliance. \emph{Training-time} defenses (StruQ~\cite{struq}, SecAlign~\cite{secalign},
Instruction Hierarchy~\cite{hierarchy}) teach a data/instruction boundary. \emph{Detection} spans
classifiers, known-answer checks~\cite{liu,datasentinel}, and task-drift probes~\cite{taskdrift}.
\emph{System-level} defenses enforce security outside the model: Dual LLM~\cite{dualllm},
CaMeL~\cite{camel}, and FIDES~\cite{fides} use capabilities or information-flow control, with strong
guarantees but utility cost. AgentShield borrows the known-answer idea for canaries and the
information-flow idea for taint, and adds the multilingual dimension none of these measure.

\section{Threat Model}
The attacker controls some retrieved content but cannot alter the user instruction, system prompt,
tools, or firewall. Goals: action hijack, data exfiltration, prompt leak, misinformation, denial.
The defender trusts the user (USER) and the user's sensitive data (PRIVATE) and treats retrieved
content (EXTERNAL) as untrusted. All evaluation uses a sandbox with mock tools and fake secrets.

\section{Method}
Every tool call passes through one hook. Incoming content is preprocessed, scored by three
detectors, fused, and then allowed, sanitized, or blocked; outgoing calls pass canary and taint
checks and per-task permissions before execution (Fig.~\ref{fig:arch}).

\begin{figure}[t]
\centering
\fbox{\parbox{0.92\columnwidth}{\footnotesize
\texttt{output} $\rightarrow$ preprocess $\rightarrow$ detectors(L1,L2,L3) $\rightarrow$ fuse $\rightarrow$ sanitize/allow/block\\[2pt]
\texttt{tool call} $\rightarrow$ canary $\rightarrow$ taint policy $\rightarrow$ permissions/approval $\rightarrow$ allow/block}}
\caption{AgentShield pipeline. Replace with a rendered diagram if desired.}
\label{fig:arch}
\end{figure}

\textbf{Preprocessing} separates visible from hidden HTML, decodes base64/hex/URL/ROT13, applies
NFKC, strips invisible characters (keeping Indic ZWNJ/ZWJ), maps homoglyphs, and OCRs images/PDFs.
\textbf{L1} is a fast multilingual rule filter with structural signals. \textbf{L2} is a fine-tuned
\texttt{xlm-roberta-base} classifier (the recall layer). \textbf{L3} is an LLM intent check
comparing the user task with the content's ask. \textbf{Canary tokens} are placed in the prompt and
sensitive files; any appearance in an outgoing call blocks and traces it. \textbf{Taint tracking}
labels data USER/PRIVATE/EXTERNAL and blocks PRIVATE$\rightarrow$EXTERNAL flows (or requires
approval). A logistic \textbf{fusion} combines layer scores with weights learned on validation data;
the \textbf{firewall} enforces the decision with a fail-safe approval timeout.

\section{Dataset}
A unified JSONL schema labels each sample by label, technique, language, script, source, and goal.
Attacks come from public corpora and a human-authored Hindi/Tamil/Hinglish/Tanglish set; benign hard
negatives (recipes, manuals, quoted text, code) stress false positives. Splits are 70/15/15,
stratified, with one attack category held out for generalization. The snapshot used here merges
public English injections with the authored negatives: test 87 (31 attack / 56 benign).

\section{Experiments}
Metrics: detection (recall), attack-success rate ($1-$recall), precision, false positives (FPR), and
latency. Baselines: no defense, a keyword filter, and an open-source classifier (scored on raw
text). The reported run uses L1 with learned fusion; L2/L3 were inactive (training/Ollama), and
canary/taint are action-level, shown qualitatively.

\section{Results}
\begin{table}[t]
\caption{Detection on the public English injection test set}
\label{tab:results}
\centering
\begin{tabular}{lcccc}
\toprule
Defense & Recall & Attack succ. & FPR & Precision \\
\midrule
No defense & 0.00 & 1.00 & 0.00 & -- \\
Keyword filter & 0.19 & 0.81 & 0.00 & 1.00 \\
\textbf{AgentShield (L1+fusion)} & \textbf{0.23} & \textbf{0.77} & \textbf{0.036} & \textbf{0.78} \\
\bottomrule
\end{tabular}
\end{table}

Table~\ref{tab:results} shows AgentShield's rule layer exceeds the keyword baseline in recall and F1
at a comparably low FPR, with latency p50 0.07\,ms / p95 0.36\,ms. Recall is bounded by L2 being
inactive; L1 is deliberately high-precision.

\textbf{Provenance result.} The central contribution is independent of detector recall: in the test
suite, a benign-looking page (no trigger words) drives the agent to email a secret, and the
canary/taint layer blocks the send and traces it to the private file --- content scanning alone
fails here. The preprocessor reverses base64/hex/URL/homoglyph obfuscations that evade the keyword
baseline.

\section{Ablation}
Removing the (only active) rule layer drops detection to zero. With L2/L3 trained/enabled, the
harness produces the full per-layer ablation; the canary/taint ablation is action-level and removes
the exfiltration backstop.

\section{Discussion and Limitations}
L1 is a precise, low-recall filter; L2 carries recall and L3 resolves ambiguity, while canary/taint
provide a deterministic backstop: to exfiltrate, attacker data must reach an external sink, which
taint/canary gate regardless of wording. Limitations: the reported numbers are L1-only on an English
snapshot; local 8B models do tool-calling imperfectly; taint is coarse for paraphrased values
(mitigated by approval).

\section{Ethics}
All attacks run against mock tools in a sandbox with fake secrets. The dataset is released for
defensive research under CC BY 4.0 with a responsible-use note. The adaptive attacker is a defensive
evaluation tool confined to the sandbox.

\section{Conclusion}
Combining multilingual detection with provenance-based canary and taint tracking yields a practical,
low-latency defense that catches exfiltration detection alone misses. The framework is complete;
training L2 and expanding the multilingual corpus scale the numbers, not the design.

\begin{thebibliography}{00}
\bibitem{greshake} K. Greshake \textit{et al.}, ``Not what you've signed up for: Compromising real-world LLM-integrated applications with indirect prompt injection,'' AISec, 2023.
\bibitem{agentdojo} E. Debenedetti \textit{et al.}, ``AgentDojo,'' NeurIPS D\&B, 2024.
\bibitem{injecagent} Q. Zhan \textit{et al.}, ``InjecAgent,'' Findings of ACL, 2024.
\bibitem{bipia} J. Yi \textit{et al.}, ``Benchmarking and defending against indirect prompt injection,'' KDD, 2025.
\bibitem{multijail} Y. Deng \textit{et al.}, ``Multilingual jailbreak challenges in LLMs,'' ICLR, 2024.
\bibitem{lowres} Z.-X. Yong \textit{et al.}, ``Low-resource languages jailbreak GPT-4,'' NeurIPS SoLaR, 2023.
\bibitem{adaptive} Q. Zhan \textit{et al.}, ``Adaptive attacks break defenses against indirect prompt injection,'' Findings of NAACL, 2025.
\bibitem{spotlight} K. Hines \textit{et al.}, ``Defending against indirect prompt injection with spotlighting,'' 2024.
\bibitem{struq} S. Chen \textit{et al.}, ``StruQ,'' USENIX Security, 2025.
\bibitem{secalign} S. Chen \textit{et al.}, ``SecAlign,'' 2024.
\bibitem{hierarchy} E. Wallace \textit{et al.}, ``The instruction hierarchy,'' 2024.
\bibitem{liu} Y. Liu \textit{et al.}, ``Formalizing and benchmarking prompt injection attacks and defenses,'' USENIX Security, 2024.
\bibitem{datasentinel} Y. Liu \textit{et al.}, ``DataSentinel,'' IEEE S\&P, 2025.
\bibitem{taskdrift} S. Abdelnabi \textit{et al.}, ``Get my drift? Catching LLM task drift with activation deltas,'' SaTML, 2025.
\bibitem{dualllm} S. Willison, ``The dual LLM pattern,'' 2023.
\bibitem{camel} E. Debenedetti \textit{et al.}, ``Defeating prompt injections by design (CaMeL),'' SaTML, 2026.
\bibitem{fides} M. Costa \textit{et al.}, ``Securing AI agents with information-flow control (FIDES),'' 2025.
\end{thebibliography}

\end{document}
